\begin{abstract}
Understanding where factual knowledge resides in transformer language models---and how to surgically remove it---is central to machine unlearning, model editing, and mechanistic interpretability. We present a unified analysis that combines logit lens profiling, component ablation, directional deletion, and geometric characterization of the \residstream in \gptmed (24 layers, 355M parameters) across 200 \counterfact examples. We find that factual information follows a \emph{non-monotonic} trajectory: it is enriched at layers 15--21 and then \emph{actively suppressed} in the final 2--3 layers. Component ablation reveals a two-stage architecture in which early MLPs (layers 0, 3, 6) serve as the primary knowledge store while late attention heads extract and route information. Projecting out a single mean fact direction from MLP outputs at layers 0--3 removes 99.7\% of the target probability, and this deletion fundamentally reshapes the \residstream geometry---cosine similarity between pre- and post-deletion activations drops to approximately zero at the final layer. Geometrically, the \residstream is dominated by a single principal component capturing 63--91\% of variance, with activation norms growing monotonically from 130 to 1055 across layers. Our results provide a unified picture of how facts are stored, amplified, suppressed, and deleted in the \residstream, with direct implications for targeted knowledge removal.
\end{abstract}
